\begin{corollary}[Lag-two studentized scalar interval]\label{cor:finite-record-lag-two-studentized-main}
Work under \Cref{ass:fixed-sampled-boundary-record-main}.  Define the observed
complete-cycle lag-two covariance estimator by
\[
\widehat\gamma_{2,N}
=\frac1N\sum_{j\in J_{2,N}}\sum_{n=\tau_j}^{\tau_{j+1}-1}
(A_n-\widehat\rho_N)(A_{n+2}-\widehat\rho_N),
\qquad
\widehat\rho_N=\frac1N\sum_{n=0}^{N-1}A_n,
\]
where \(J_{2,N}\) contains exactly those all-complete-cycle indices for which
the lag-two windows lie in the observed record.
This is the paper's fully studentized scalar visible-record interval.

The scalar studentizer used for this statistic is the following observable
separated-batch estimator.  Let
\[
b_N=\lfloor N^\beta\rfloor,\qquad 0<\beta<1,\qquad d_N=b_N+1,
\]
and let
\[
M_{2,N}=\max\{m\ge0:(m-1)d_N+b_N-1\in J_{2,N}\},
\qquad
I^{\gamma}_{\ell,N}
=\{(\ell-1)d_N,\ldots,(\ell-1)d_N+b_N-1\},
\]
with \(M_{2,N}=0\) if the displayed set is empty.  For
\(j\in J_{2,N}\), define the plug-in lag-two cycle reward
\begin{equation}\label{eq:lag-two-plugin-cycle-reward-main}
\widehat Y^{\gamma}_{j,N}
=
\sum_{n=\tau_j}^{\tau_{j+1}-1}
\left\{(A_n-\widehat\rho_N)(A_{n+2}-\widehat\rho_N)
-\widehat\gamma_{2,N}\right\}.
\end{equation}
Put
\[
\widehat C^{\gamma}_{\ell,N}
=\sum_{j\in I^{\gamma}_{\ell,N}}\widehat Y^{\gamma}_{j,N},
\qquad
\bar{\widehat C}^{\gamma}_N
=M_{2,N}^{-1}\sum_{\ell=1}^{M_{2,N}}\widehat C^{\gamma}_{\ell,N}
\quad\text{on }\{M_{2,N}\ge1\}.
\]
The scalar lag-two variance estimator is reported only when at least two
complete separated lag-two batches are available.  For the proof of consistency
on a total probability space, use the auxiliary completed statistic
\begin{equation}\label{eq:lag-two-scalar-studentizer-main}
\widehat\sigma^{2,\rm aux}_{\gamma_2,N}
=
\begin{cases}
\displaystyle
\frac{K_N}{N}\,
\frac1{b_N(M_{2,N}-1)}
\sum_{\ell=1}^{M_{2,N}}
\left(\widehat C^{\gamma}_{\ell,N}
-\bar{\widehat C}^{\gamma}_N\right)^2,
&M_{2,N}\ge2,\\[1.1ex]
0,&M_{2,N}<2 .
\end{cases}
\end{equation}
On \(\{M_{2,N}<2\}\) the lag-two interval is not reported; the zero branch in
\(\widehat\sigma^{2,\rm aux}_{\gamma_2,N}\) is used only in the convergence
proof.  On \(\{M_{2,N}\ge2\}\), put
\(\widehat\sigma^2_{\gamma_2,N}:=
\widehat\sigma^{2,\rm aux}_{\gamma_2,N}\).
Then
\[
\sqrt N(\widehat\gamma_{2,N}-\gamma_2)
\Rightarrow \mathcal N(0,\sigma_{\gamma_2}^2),
\]
where \(\sigma_{\gamma_2}^2\) is the explicit positive variance
\eqref{eq:lag-two-explicit-variance-main}, and
\[
\widehat\sigma^{2,\rm aux}_{\gamma_2,N}
\xrightarrow{\PP}\sigma_{\gamma_2}^2 .
\]
Consequently, for \(0<\alpha<1\), with
\(z_{1-\alpha/2}\) denoting the standard-normal \(1-\alpha/2\) quantile and
\(\widehat\sigma_{\gamma_2,N}=(\widehat\sigma^2_{\gamma_2,N})^{1/2}\),
the fixed-setting two-sided interval reported on \(\{M_{2,N}\ge2\}\),
\begin{equation}\label{eq:lag-two-studentized-interval-main}
\left[
\widehat\gamma_{2,N}
-z_{1-\alpha/2}\frac{\widehat\sigma_{\gamma_2,N}}{\sqrt N},
\quad
\widehat\gamma_{2,N}
+z_{1-\alpha/2}\frac{\widehat\sigma_{\gamma_2,N}}{\sqrt N}
\right]
\end{equation}
has limiting coverage \(1-\alpha\).
\end{corollary}

\begin{proof}
\Cref{lem:lag-two-reward-main} identifies the complete-cycle reward for the
true-centered statistic as a centered \(1\)-dependent reward, and
\Cref{lem:all-cycle-lag-two-windows-main} verifies that every lag-two term
used by \(J_{2,N}\) is observed.  Thus the \(m_{\rm dep}=1\) case of
\Cref{lem:finite-lag-renewal-reward-main} applies with
\[
\sigma_{\gamma_2}^2
=\frac{\Var(Y_0)+2\operatorname{Cov}(Y_0,Y_1)}{\EE(G+1)}.
\]
\Cref{lem:lag-two-variance-nondegenerate-main} identifies this quantity with
the double/triple sum \eqref{eq:lag-two-explicit-variance-main} and proves
that it is strictly positive on the full positive parameter domain.
Replacing \(\rho\) by \(\widehat\rho_N\) in the displayed estimator changes
the true-centered average by
\[
-(\widehat\rho_N-\rho)\frac1N
\sum_{j\in J_{2,N}}\sum_{n=\tau_j}^{\tau_{j+1}-1}
\bigl((A_n-\rho)+(A_{n+2}-\rho)\bigr)
+(\widehat\rho_N-\rho)^2\frac1N\sum_{j\in J_{2,N}}L_j .
\]
The stationarity and exponential renewal tail give the ordinary record mean
\(\widehat\rho_N-\rho=O_{\PP}(N^{-1/2})\).  The trimmed bracket is exactly the
linear-combination estimate
\eqref{eq:trimmed-shifted-linear-combination-main} with
\(\Lambda=\{0,2\}\) and \(c_0=c_2=1\):
\[
\frac1N
\sum_{j\in J_{2,N}}\sum_{n=\tau_j}^{\tau_{j+1}-1}
\bigl((A_n-\rho)+(A_{n+2}-\rho)\bigr)
=O_{\PP}(N^{-1/2}).
\]
Also \(N^{-1}\sum_{j\in J_{2,N}}L_j=O_{\PP}(1)\).  Therefore the two plug-in
terms are respectively \(O_{\PP}(N^{-1})\) and \(O_{\PP}(N^{-1})\), hence
\(o_{\PP}(N^{-1/2})\).  The CLT therefore has the same variance as the
true-centered reward.

It remains only to identify the displayed scalar studentizer with the
batch-estimated variance of that same limiting reward.  Let
\[
Y^\gamma_j
=\sum_{n=\tau_j}^{\tau_{j+1}-1}
\{(A_n-\rho)(A_{n+2}-\rho)-\gamma_2\},
\qquad
C^\gamma_{\ell,N}=\sum_{j\in I^\gamma_{\ell,N}}Y^\gamma_j .
\]
Write
\(\bar C^\gamma_N=M_{2,N}^{-1}\sum_{\ell=1}^{M_{2,N}}C^\gamma_{\ell,N}\)
on \(\{M_{2,N}\ge1\}\).
The \(m_{\rm dep}=1\) case of
\Cref{lem:finite-lag-separated-batch-main} applied to
\((Y^\gamma_j)\) gives consistency of
\[
\frac{K_N}{N}\,
\frac1{b_N(M_{2,N}-1)}
\sum_{\ell=1}^{M_{2,N}}
\left(C^\gamma_{\ell,N}-\bar C^\gamma_N\right)^2
\]
for \(\sigma_{\gamma_2}^2\), with the same auxiliary zero completion on
\(\{M_{2,N}<2\}\), an event whose probability tends to zero and on which the
lag-two interval is not reported.

For the plug-in batch rewards, write
\[
\Delta_{\rho,N}=\widehat\rho_N-\rho,\qquad
\Delta_{\gamma,N}=\widehat\gamma_{2,N}-\gamma_2,
\]
and
\[
W_j=\sum_{n=\tau_j}^{\tau_{j+1}-1}
\bigl\{(A_n-\rho)+(A_{n+2}-\rho)\bigr\}.
\]
Then, for every used cycle,
\begin{equation}\label{eq:lag-two-plugin-cycle-difference-main}
\widehat Y^\gamma_{j,N}-Y^\gamma_j
=-\Delta_{\rho,N}W_j
+(\Delta_{\rho,N}^2-\Delta_{\gamma,N})L_j .
\end{equation}
The CLT already proved above gives
\(\Delta_{\gamma,N}=O_{\PP}(N^{-1/2})\), and the ordinary renewal mean gives
\(\Delta_{\rho,N}=O_{\PP}(N^{-1/2})\).  Since \(|W_j|\le2L_j+4\),
\Cref{prop:calendar-stopped-separated-batch-plugin-main} gives the same uniform
\(O_{\PP}(b_NN^{-1/2})\) block perturbation as
\eqref{eq:separated-batch-plugin-block-main}.  The plug-in covariance
perturbation bound \eqref{eq:separated-batch-plugin-covariance-main}, applied
in this scalar \(m=1\) setting, is \(o_{\PP}(1)\) because \(b_N/N\to0\).
Thus
\(\widehat\sigma^{2,\rm aux}_{\gamma_2,N}\to_{\PP}\sigma_{\gamma_2}^2\), and
the reported \(\widehat\sigma^2_{\gamma_2,N}\) has the same limit on the
retained event.
Since the limiting variance is positive by
\Cref{lem:lag-two-variance-nondegenerate-main}, Slutsky's theorem gives the
studentized statistic and hence the interval
\eqref{eq:lag-two-studentized-interval-main}.
\end{proof}

\begin{corollary}[Side-data-only labelled-rate transport after the quotient inverse]\label{cor:side-data-labelled-rate-wald-main}
This corollary is not a visible-record parameter-estimation result.  It is the
ordinary delta-method transport obtained after the quotient-only estimator
\(\Phi(\Rinch_N)\) has been formed and after external side data have
chosen an ordering of the two separated roots.  It also conditions on the
sampling interval \(\Delta\tau\) being a known design interval; otherwise the
same construction transports only the ordered dimensionless products.  Work
under the fixed finite-record experiment and reporting gate \(\mathcal E_N\).
Suppose the quotient population point is separated:
\[
\Delta_\sigma=(p-s)^2>0,\qquad p,s\in(0,1),
\]
and suppose the visible-record estimator satisfies the inclusion CLT and
covariance rate used in \Cref{thm:finite-record-inference-audit-main}:
\[
\sqrt N(\Rinch_N-\Rinc)\Rightarrow \mathcal N(0,\Sigma_r),\qquad
\|\widehat\Sigma_{r,N}^{\rm tot}-\Sigma_r\|_{\operatorname{op}}
=o_{\PP}(t_N)
\]
for the retained-rank threshold \(t_N\).  Add one of the following side-data
hypotheses, neither of which is a consequence of the visible record:

\begin{enumerate}[label=\textup{(\roman*)},leftmargin=2.2em,itemsep=0.35em]
\item A deterministic oracle branch map \(\mathcal B\) is fixed on a
neighbourhood of the unordered root set and orders the two roots as the
attempt and recovery branches.
\item A calibration or side-data estimator \(\widehat{\mathcal B}_N\) is used
only on an event \(\mathcal C_N\) such that
\[
\PP(\mathcal C_N)\to1,\qquad
\mathcal C_N\subseteq
\{\widehat{\mathcal B}_N=\mathcal B
\text{ on the retained separated-root neighbourhood}\}.
\]
\end{enumerate}

Then the visible-record part of the construction supplies only
\(\Phi(\Rinch_N)\) and the unordered root set.  On \(\mathcal E_N\) and on
the corresponding side-data retained chart, after applying the external branch
map, the ordered labelled-rate transport
\[
(\widehat\Gamma_N,\widehat\kappa_{\mathrm r,N})
=\left(-\frac{\log\widehat p_N}{\Delta\tau},
       -\frac{\log\widehat s_N}{\Delta\tau}\right)
\]
satisfies
\[
\sqrt N\{(\widehat\Gamma_N,\widehat\kappa_{\mathrm r,N})
-(\Gamma,\kappa_{\mathrm r})\}
\Rightarrow
\mathcal N(0,\Sigma_{\Gamma,\kappa}),
\qquad
\Sigma_{\Gamma,\kappa}
=J_{\Gamma,\kappa}(\Rinc)\Sigma_rJ_{\Gamma,\kappa}(\Rinc)^\top .
\]
With
\[
\widehat\Sigma_{\Gamma,\kappa,N}
=\widehat J_{\Gamma,\kappa,N}\widehat\Sigma_{r,N}
\widehat J_{\Gamma,\kappa,N}^{\top},
\]
where \(\widehat J_{\Gamma,\kappa,N}\) is
\eqref{eq:inclusion-to-rate-jacobian-main} evaluated on the retained labelled
chart, the covariance is consistent at the same operator-norm scale.  If
\[
\|\widehat\Sigma_{\Gamma,\kappa,N}
-\Sigma_{\Gamma,\kappa}\|_{\operatorname{op}}=o_{\PP}(t_N),
\]
then the thresholded retained-rank Wald statistic converges to
\(\chi_k^2\), \(k=\operatorname{rank}\Sigma_{\Gamma,\kappa}\).  No labelled-rate
conclusion is asserted on \(\Delta_\sigma=0\), without an oracle map, or for an
estimated branch map outside the high-probability agreement event.  In
particular, this corollary never converts Theorem~B into identification of an
ordered hidden labelled D-MAP representation from the visible law alone.
\end{corollary}

\begin{proof}
On a separated-root neighbourhood the oracle branch map fixes one \(C^1\)
ordering of the two roots, and
\Cref{lem:separated-root-rate-delta-main} gives the root and rate Jacobians.
The complete-cycle reporting gate has probability tending to one at the stable
population point, and the separated discriminant/end-point thresholds are
eventually retained.  The multivariate delta method therefore gives the
displayed Gaussian limit after the external branch map has supplied the
ordering.  If an estimated map is used, the event
\(\mathcal C_N\) makes the estimated-map estimator identical to the oracle-map
estimator with probability tending to one, so the same limit follows by the
standard high-probability replacement argument.  The covariance consistency is
the same continuous-mapping argument applied to the plug-in Jacobian on the
retained separated chart.  Finally, Weyl's eigenvalue inequality retains
exactly the positive limiting covariance eigenvalues above the threshold
\(t_N\), and Slutsky's theorem gives the retained-rank \(\chi_k^2\) limit.
Every step after \(\Phi(\Rinch_N)\) is conditional on side-data ordering;
the quotient estimator and its covariance remain the visible-record output.
\end{proof}

\begin{proposition}[Quotient-only estimands from one constrained D-MAP fixed-setting long record]\label{thm:one-record-estimates-main}
Under \Cref{ass:fixed-sampled-boundary-record-main}, a single fixed-setting
record estimates only quantities of the stationary visible law: the inclusion
coordinates, the projected
symmetric coordinates \((p+s,ps)\), the unordered survival-factor pair and
unordered numerical rate set when \(\Delta\tau\) is a known design interval,
the corresponding dimensionless product set otherwise, the fixed-setting
lag-two covariance, and their regenerative covariances under the assumptions
above.  It does not label two separated rates and does not identify any object
outside the visible stationary record law.  Response loci, KMS quantities,
spectral weights, shell locations, ordered hidden labels, and ordered
\((\Gamma,\kappa_{\mathrm r})\) are therefore not outputs of the one-record
CLT or of the quotient inverse.  The branch-map convention is
\Cref{conv:branch-map-labelled-rates-main}; at \(\sigma=(2z,z^2)\) the common
dimensionless product, and only for known \(\Delta\tau\) the common numerical
rate, is the invariant visible-law exception described in
\Cref{lem:diagonal-common-rate-main}.
\end{proposition}

\begin{proof}
The positive estimands are exactly the CLT targets in
\Cref{thm:finite-record-inference-audit-main,cor:finite-record-lag-two-studentized-main}
together with the quotient transform in
\Cref{thm:toeplitz_image_inverse_duality,thm:finite-record-quotient-delta-main}.
The non-identifiability statement is
\Cref{prop:hidden_label_exchange_nonidentifiability}: exchanging the two
hidden labels leaves every finite-dimensional visible distribution unchanged.
Hence a statistic of one visible record can be a function only of the quotient
law, not of a chosen hidden ordering.  Side-data branch maps enter only through
\Cref{conv:branch-map-labelled-rates-main,cor:side-data-labelled-rate-wald-main}.
Response-locus statements are appendix/supplementary conditional corollaries
that require an independently calibrated response closure beyond the
visible-record law.  The record analysis starts from \(T_0,T_1\) and therefore
supplies no estimator for
any externally defined scalar response map, shell location, KMS parameter, or
inverse image of such a map.
\end{proof}
